\documentclass[10pt,a4paper]{amsart}
\usepackage{amsmath,amssymb,amsthm,graphicx}
\usepackage{hyperref}
\usepackage[capitalize, nameinlink]{cleveref}
\crefname{theorem}{Theorem}{Theorems}
\crefname{proposition}{Proposition}{Propositions}
\crefname{observation}{Observation}{Observations}
\crefname{lemma}{Lemma}{Lemmas}
\crefname{claim}{Claim}{Claims}
\crefname{problem}{Problem}{Problems}
\crefname{conjecture}{Conjecture}{Conjectures}
\crefname{question}{Question}{Questions}
\crefname{example}{Example}{Examples}
\crefname{fact}{Fact}{Facts}
\usepackage{amssymb}
\newcommand{\N}{\mathbb{N}}
\renewcommand{\restriction}[2]{#1 \upharpoonright #2}
\newtheorem{theorem}{Theorem}[section]
\newtheorem{lemma}[theorem]{Lemma}
\newtheorem{Claim}[theorem]{Claim}
\newtheorem{claim}[theorem]{Claim}
\newtheorem{subclaim}[theorem]{Subclaim}
\newtheorem{corollary}[theorem]{Corollary}
\newtheorem{proposition}[theorem]{Proposition}
\newtheorem{conjecture}[theorem]{Conjecture}
\newtheorem{question}[theorem]{Question}
\newtheorem{problem}[theorem]{Problem}
\newtheorem{fact}[theorem]{Fact}
\newtheorem{ax}{Axiom}


\theoremstyle{definition}

\newtheorem{definition}[theorem]{Definition}
\newtheorem{example}[theorem]{Example}
\newtheorem{remark}[theorem]{Remark}
\DeclareMathOperator{\graph}{graph}
\DeclareMathOperator{\dist}{dist}
\DeclareMathOperator{\didistance}{didist}
\newcommand{\mc}{\mathcal}
\newcommand{\rstr}{\restriction}
\newcommand{\vfi}{\varphi}
\newcommand{\sse}{\subseteq}
\newcommand{\est}{\emptyset}


\newcommand{\treegraph}[1]{Tree(#1)}
\newcommand{\fG}{\mathcal{G}}
\newcommand{\fA}{\mathcal{A}}
\newcommand{\fB}{\mathcal{B}}
\newcommand{\fC}{\mathcal{C}}
\newcommand{\fD}{\mathcal{D}}
\newcommand{\fE}{\mathcal{E}}
\newcommand{\Epsilon}{\fE}
\newcommand{\fF}{\mathcal{F}}
%\newcommand{\fG}{\mathcal{G}}
\newcommand{\fH}{\mathcal{H}}
\newcommand{\fI}{\mathcal{I}}
\newcommand{\fJ}{\mathcal{J}}
\newcommand{\fK}{\mathcal{K}}
\newcommand{\fL}{\mathcal{L}}
\newcommand{\fM}{\mathcal{M}}
\newcommand{\fN}{\mathcal{N}}
\newcommand{\fO}{\mathcal{O}}
\newcommand{\fP}{\mathcal{P}}
\newcommand{\fQ}{\mathcal{Q}}
\newcommand{\fR}{\mathcal{R}}
\newcommand{\fS}{\mathcal{S}}
\newcommand{\fT}{\mathcal{T}}
\newcommand{\fU}{\mathcal{U}}
\newcommand{\fV}{\mathcal{V}}
\newcommand{\fW}{\mathcal{W}}
\newcommand{\fX}{\mathcal{X}}
\newcommand{\fY}{\mathcal{Y}}
\newcommand{\fZ}{\mathcal{Z}}
\title{Hyperfiniteness of a countable Borel equivalence relation on a positive Ramsey cone}
\author{Balázs Bursics, Zoltán Vidnyánszky}
\date{arxiv:2412.01315v2; CC BY 4.0}
\begin{document}
\maketitle
\noindent\textit{Editorial note.} The author explicitly assumes metric closedness for the Ramsey space at470, so the counted statement is read with that source-global convention. Only the listed classical Ramsey-space properties are free inputs; the separated-cover criterion and fusion scheduling are proved.\par\medskip
\noindent\textit{Editorial note.} In641 the indices of theB-set interval are shifted by one stage. The construction638-639 assigns the cones at stage n to k\_n<k<=k\_(n+1). A point C in[A] is below A\_(n+1), so its relevant prefix cone lies in the B\_k selected at that same stage. The original641 is preserved and this bookkeeping interpretation requires independent confirmation.\par\medskip
\noindent\textit{Editorial note.} In the diameter argument602-605, a component with no current marker is unchanged. A component with one marker consists of its old component together with the old components joined by the new edges of G\_(n+1) restricted to B; all paths then have length at most2(f(n)+1). The displayed union604 is read within that restricted vertex set.\par\medskip
\noindent\textit{Editorial note.} The later generic-comeager theorem664-676 is not used or counted here; no inference about saturation of a meager set is needed for this task.\par\medskip
\noindent\textit{Editorial scope.} One arbitrary-topological-Ramsey-space hyperfiniteness theorem, with the full separated-cover and fusion proof chains. Its Ellentuck special case and the old bounded-degree lemma are not counted or used as a shortcut. Standard Borel equivalence and abstract Ramsey facts are declared precisely. Source index/union conventions remain explicitly subject to independent review.\par\medskip
\noindent\textit{Allowed standard prerequisites (editorial).}\begin{itemize}
\item Classical topological Ramsey-space fusion and Ramsey properties. Use the metric-closed topological Ramsey-space setting stipulated in source470. Basic cones[a,A] are Borel, each finite downset\{s:s<=fin r\_n(A)\} is finite, fusion sequences have the stated limit, and the two fusion conclusions474-480 and finite Borel-color homogenization484-486 hold.. Todorcevic, Introduction to Ramsey Spaces; exact properties supplied in source449-486.. Provides individual homogeneous refinements and fusion limits, not the separated-cover scheduling or its hyperfinite finite-component construction.
\item Classical countable Borel equivalence-relation facts. A CBER is the union of countably many graphs of Borel involutions. Hyperfiniteness is equivalent to being an increasing union of finite Borel equivalence relations, and is inherited by restriction to a Borel subset.. Convert the constructed finite components into hyperfiniteness and represent an arbitrary E by bounded-degree graphs.
\item Finite Borel vertex coloring for bounded-degree irreflexive graphs. Such a graph has a finite Borel proper vertex coloring (Kechris-Solecki-Todorcevic).. Source625.. Applied to a bounded graph-distance power before the Ramsey homogeneous refinement.
\end{itemize}\medskip
	A \emph{countable Borel equivalence relation (CBER)} on a Polish space $X$ is an equivalence relation with all classes countable that is Borel as a subset of $X^2$. CBERs have a rich and well investigated structure theory (see, e.g., \cite{kechris2024theory}). The natural way to compare two such equivalence relations is the notion of \emph{reduction}: if $E$ and $F$ are CBERs on spaces $X$ and $Y$, $E$ is said to \emph{Borel reduce} to $F$ if there exists a Borel map $\phi:X \to Y$ with
	\[\forall x,x' \in X \ xEx' \iff \phi(x)F\phi(x').\]
	
	The simplest non-trivial class of CBERs is the collection of \emph{hyperfinite} ones, that is, the collection of CBERs which admit a Borel reduction to $\mathbb{E}_0$, the eventual equality equivalence relation on $2^\N$. 

	Recall the following definition. 
	\begin{definition}
		Let $G$ be a graph, $B \subseteq V(G)$. The set $B$ is called \emph{$k$-separated} if for any $x \neq x' \in B$ we have $\dist_{G}(x,x')>k$.
	\end{definition}

	A \emph{topological Ramsey space} in the sense of Todor\v{c}evi\'c \cite{ramseyspaces} is a triple $(\mathcal{R},\le,r)$ satisfying a certain set of axioms, where $\mathcal{R}$ is a nonempty set, $\le$ a quasi-ordering on $\mathcal{R},$ and $r:\mathcal{R}\times\N\to\mathcal{AR}$. The range of $r$ can be thought of as the collection of finite approximations to elements of $\mathcal{R}$. We will use capitals $A,B,\ldots$ for the elements of $\mathcal{R}$, and $a,b,\ldots$ for their approximations. We denote $r_n(.)=r(.,n)$.
	
	The relation $\le$ is encoded on the level of approximations by a quasi-ordering $\le_{\text{fin}}$ on $\mathcal{AR}$.
	
	For $a\in\mathcal{AR}$ and $B\in\mathcal{R}$ define 
	$$[a,B]=\{A\in\mathcal{R}:A\le B \wedge (\exists n)\; r_n(A)=a\}.$$
	
	We also use the abbreviations $[B]=\{A\in\mathcal{R}:A\le B \}$ and $[n,B]=[r_n(B),B].$
	
	There are multiple natural topologies on $\fR$: the \emph{metric topology} of $\fR$ inherited from the Baire space $\fA\fR^\N,$ and the so-called \emph{Ellentuck topology} generated by $\{[a,A]:\; a\in\fA\fR,\; A\in\fR\}.$ We work with the metric topology, unless stated otherwise.
	
	
	
	The next definition gives a standard way of constructing elements of $\fR$ with some desirable property.
	
	\begin{definition}
		A \emph{fusion sequence} is a sequence $(A_n)_{n\in \N}$ in $\fR$ such that $A_{n+1}\in [n,A_n]$ for every $n\in\N.$
		
		The \emph{limit} of a fusion sequence $(A_n)_{n\in\N}$, denoted $\lim A_n$, is the unique $A\in\fR$ such that $A\in[n,A_n]$ for every $n\in\N.$ 
	\end{definition}
	
	Note that for the existence of the limit of a fusion sequence we assume that $\fR$ is closed in $\fA\fR^\N$ with respect to the metric topology. This is not listed among the axioms of topological Ramsey spaces in \cite{ramseyspaces}, however, this property is assumed for the main results of the theory. We make use of the following properties of fusion sequences.


	
	\begin{proposition}\label{fusion}
		Let $(A_n)$ be a fusion sequence and $C\in[\lim A_n].$ Then 
        \begin{enumerate}
            \item $\lim A_n\le A_k$ for every $k$,
            \item there are infinitely many $n$ with $C\in[s,A_n]$ for some $s\le_{fin}r_n(A_n).$
        \end{enumerate}
	\end{proposition}
	
	Apart from the technique of fusion sequences, we use the Ramsey theoretic statement below.
	
	\begin{theorem}\label{ramsey property}
		Let $\fR$ be a topological Ramsey space, $c:\fR\to k$ a finite Borel coloring, $A\in\fR,$ and $s\le_{fin}r_n(A)$. Then there exists $B\in[n,A]$ such that $\rstr{c}{[s,B]}$ is monochromatic.
	\end{theorem}

	\begin{theorem} \label{t:ramsey hyperfin}
		Let $\mathcal{R}$ be a topological Ramsey space and $E$ be a CBER on $\mathcal{R}$. Then there is a set $A \in \mc{R}$ such that $\restriction{E}{[A]}$ is hyperfinite. 
	\end{theorem}

	\begin{theorem}
		\label{t:separation2}
		Let $E$ be a CBER on the space $X$, let $(G_n)_{n \in \N}$ be an increasing sequence of bounded degree Borel graphs\footnote{Throughout the paper, graphs are not necessarily assumed to be irreflexive.} such that $\bigcup_n G_n=E$, and let $f:\N\to\N$ be such that $\forall n\in\N\; f(n+1)\ge2\cdot(f(n)+1)$. Moreover, assume that $B_n \subseteq X$ are Borel sets so that every $B_n$ is $f(n)$-separated in $G_n$. 
		
		Then $\restriction{E}{B}$ is hyperfinite, where $B=\{x:\exists^\infty n \ (x \in B_n)\}$.
	\end{theorem}

	\begin{proof}
		
		We define a sequence $H_0 \subseteq H_1 \subseteq \dots \subseteq H_n \subseteq \dots \subseteq E$ inductively. Let $H_0=\emptyset$. 
		
		Given $H_n$, let \[H_{n+1}=H_{n} \cup \{(x,y) \in \restriction{G_{n+1}}{B}:\; x\in B_{n+1} \text{ or } y\in B_{n+1}\}.\]
		
		First we check that $\rstr{\bigcup_{n \in \N} H_n}{B}=\restriction{E}{B}$. Indeed, let $(x,y) \in \rstr{G_n}{B}$. Then there is some $k \geq n$ such that $x \in B_k$. But then $(x,y) \in H_k$. 
		
		\begin{claim}
			Every connected component of $H_n$ has diameter $\le f(n)$. 
			
		\end{claim}
		\begin{proof}
			We prove this by induction on $n$. The case $n=0$ is clear, so assume that we have shown the statement for $n$. We claim that if $H$ is a connected component of $H_{n+1}$ then it contains at most one point of $B_{n+1}$. Indeed, otherwise let $x \neq x' $ from $B_{n+1}$ in $V(H)$ be such that for $p=(x_0,\dots,x_k)$, the injective path with minimal length from $x_0=x$ to $x_k=x'$ in $H$ we have $x_1,\ldots,x_{k-1}\notin B_{n+1}$. Then $x_1,\ldots,x_{k-1}$ are contained in a single connected component $H'$ of $H_n$, and by the inductive hypothesis $\dist_{H_n}(x_1,x_{k-1})\le f(n).$ But this implies $\dist_{G_{n+1}}(x,x')\le f(n)+2\le f(n+1),$ a contradiction.
			
			Thus, $H$ contains at most one point, say $x$, from $B_{n+1}$. Then \[V(H)=\bigcup \{[y]_{H_{n}}:(x,y) \in G_{n+1}\}.\]
			yielding \[\text{diam}(H) \le 2\cdot (f(n)+1) \le f(n+1).\]
		\end{proof}
		
		As $H_n\sse G_n$ which has bounded degrees, this, in turn also ends the proof of Theorem \ref{t:separation2}.
	\end{proof}

	\begin{lemma}
		\label{l:graph canonization}
		Let $\fR$ be a topological Ramsey space, let $G$ be a bounded degree Borel graph on $\fR$, let $k\in \N$, and let $a\in \fA\fR$ and $A\in\fR$ be such that $a\le_{\text{fin}}r_n(A)$. Then there exists $A'\in [n,A]$ such that $[a,A']$ is $k$-separated.
	\end{lemma}
	
	
	
	\begin{proof}
		By \cite{KST}, every irreflexive bounded degree Borel graph admits a finite Borel vertex coloring. Let $c:\fR \to m$ be such a coloring of \[G^k=\{(x,y):x\neq y \wedge \dist(x,y)\le k \}.\] 
        By Theorem \ref{ramsey property}. there is some $A'\in[n,A]$ such that $\rstr{c}{[a,A']}$ is constant, thus, $[a,A']$ is $k$-separated. 
	\end{proof}
	
	
	\begin{lemma}\label{ramsey separation}
		Let $\fR$ be a topological Ramsey-space, $G_n$ a sequence of bounded degree Borel graphs on $\fR$, and $f:\N\to\N$. Then there exist $A\in\fR$ and $B_n\sse\fR$ Borel sets such that every $B_n$ is $f(n)$-separated in $G_n$, and 
		$$[A]\sse\{C:\exists^\infty n \ (C \in B_n)\}.$$
	\end{lemma}
	
	\begin{proof}
		We construct a fusion sequence $(A_n)_{n\in\N}$ inductively. Let 
		\[k_n=\sum_{i<n}|\{s\in\fA\fR: \; s\le_{\text{fin}}r_i(A_i)\} |.\]
		Let $B_0=\est.$ Throughout the construction, we define $B_{k_n+1}, \ldots, B_{k_{n+1}}$ after fixing $A_n$ (and  $k_{n+1}$).
		Suppose that $A_n$ and $B_0,\ldots, B_{k_n}$ are given. We enumerate $\{s\in\fA\fR: \; s\le_{\text{fin}}r_n(A_n)\}$ as $\{s_j: 1\le j\le k_{n+1}-k_n\}$, and define $(A_j')_{j\le k_{n+1}-k_n}$ recursively as follows: let $A_0'=A_n$, and if $A_j'$ is already given for some $j<k_{n+1}-k_n$, choose $A_{j+1}'\in [n,A_j']$ so that $[s_{j+1},A_{j+1}']$ is $f(k_n+j+1)$-separated in $G_{k_n+j+1}$ (such an element of $\fR$ exists by Lemma \ref{l:graph canonization}), and let $B_{k_n+j+1}=[s_{j+1},A_{j+1}']$. Afterwards, let $A_{n+1}=A_{k_{n+1}-k_n}'\in [n,A_n].$
		
		Finally, let $A=\lim A_n$ and suppose that $C'\in [A].$ By Proposition \ref{fusion}. there are infinitely many $n$ with $C'\in[s,A_n]$ for some $s\le_{fin}r_n(A_n)$. For such an $n$ and $s$, we have that $[s,A_n]\sse B_k$ for some $k_{n-1}<k\le k_n,$ which yields $C' \in\{C:\exists^\infty n \ (C \in B_n)\}$.
	\end{proof}

	\begin{proof}
		As all classes of $E$ are countable, there are Borel involutions $\vfi_n$ on $\fR$ such that $E=\bigcup_{n\in\N}\graph(\vfi_n)$. Let $G_n=\bigcup_{i\in n}\graph(\vfi_i)$, and $f:\N\to \N$ be as in Theorem \ref{t:separation2}. By Lemma \ref{ramsey separation}. there exist $A\in\fR$ and $B_n\sse\fR$ Borel sets such that every $B_n$ is $f(n)$-separated in $G_n$, and 
		$$[A]\sse\{C:\exists^\infty n \ (C \in B_n)\},$$
		thus, $\rstr{E}{[A]}$ is hyperfinite according to Theorem \ref{t:separation2}.
	\end{proof}

\bibliographystyle{plain}
\bibliography{ref}
\end{document}
